% !PRINTED 50
\weberpage{50}
%<<CONT
group of the $B$, which is thus also a divisor of $S$, we shall
denote by $U$ and call the \gesperrt{group reciprocal to $T$}. Here too,
however, it is to be remarked that the concept of the reciprocal
group depends in general on the choice of the basis and of the
roots of unity $\varepsilon$, and thus does not belong to the
groups $S$ and $T$ as such.

The group reciprocal to $T'$ is characterised by the following
theorem:
% [sic: Weber prints T' here; the newly defined group is reciprocal to T.]

8.\ \gesperrt{If one lets $A$ run through the elements of a
divisor $T$ of $S$ and seeks all elements $B$ of $S$ which for
every $A$ satisfy the condition}
\weq{7}{\chi_B(A) = 1}
\gesperrt{then $B$ runs through the elements of the group $U$
reciprocal to $T$, whose degree is equal to the index of $T$.}

By the theorem of \S.\,11, (20), $T'$ is the reciprocal group to
$U$, and the relation of these two groups is therefore mutual.
% [sic: Weber prints T' here; reciprocity requires T.]

Of these reciprocal groups there holds further the theorem:

9.\ \gesperrt{If $T'$ is a divisor of $T$, then conversely the
reciprocal group $U$ of $T'$ is a divisor of the group $U'$
reciprocal to $T''$.}
% [sic: Weber prints T' and T'' here; the theorem requires T and T', respectively.]

For every element $A'$ of $T''$ is at the same time contained in
$T'$, and consequently, if $B$ runs through the group $U$,
% [sic: Weber prints T'' and T' here; the proof requires T' and T, respectively.]
$\chi_B(A') = 1$. Hence $B$ must occur in the group $U'$.

If from the totality of the characters $\chi$ one selects an
arbitrary number $\xi_1, \xi_2 \dots \xi_\mu$, no matter whether
they form a group or not, then all elements $A$ which satisfy
the $\mu$ conditions
\weq{8}{\xi_1(A) = 1,\ \xi_2(A) = 1, \dots \xi_\mu(A) = 1}
form a group, because from $\xi_i(A) = 1$, $\xi_i(A') = 1$ it
follows that $\xi_i(AA') = 1$ also. If the
$\xi_1, \xi_2 \dots \xi_\mu$ are not groups, then from the
equations (8) there follow so many further
$\xi_{\mu+1}(A) = 1 \dots$ that the
$\xi_1, \xi_2 \dots \xi_\mu$ are completed to a group.

10.\ \gesperrt{From Theorem 7. it follows that in this way all
possible divisors of $S$ can be obtained.}
